\newcommand{\ProteinDataCaption}{Original graph-digitized Tang et al. (2018) 100 mL isothermal lysozyme observations at 20 degrees C. (a) Dissolved concentration with assumed 5\% standard deviations; the diamond fixes the initial condition, filled circles train and open squares test. (b) Training conditional number densities at 74, 180 and 420 min. (c) Held-out 660-min density. (d) Conditional number $d_{10},d_{50},d_{90}$; connecting dotted lines are guides only. All sizes and quantiles condition on 2.5--25 $\mu$m. Dashed density guides use a fixed 0.6 micrometre Gaussian display bandwidth, reflected at the conditioning window and unit-area normalized. Symbols retain the original readings; the likelihood and supplied CDFs are unchanged. No model prediction or invented PSD uncertainty is shown.}
\newcommand{\ProteinFitDetails}{(a) Concentration; (b) training conditional number densities at 74, 180 and 420 min; (c) held-out 660-min density; (d) conditional number quantiles. All densities and quantiles condition on 2.5--25 $\mu$m. Solid curves show the inference-coordinate MAP; markers retain the graph readings and dashed density guides have a fixed 0.6 micrometre display bandwidth, without changing likelihood observations. Shading is the pointwise 95\% latent interval from screened NUTS draws, conditional on fixed observation errors and physical constants; it excludes additional measurement noise. The vertical dotted line marks the end of training at 420 min. The 660-min PSD and later concentrations are predicted without refitting. Quantile segments are omitted when the conditioning window is empty; quantiles derived from PSDs are not independent likelihood data. Numerical density curves and bands connect the stored bin averages at their bin centers; no smoothing or extra forward solves are used.}
\newcommand{\ProteinCornerDetails}{Diagonals show normalized posterior KDEs and dashed analytical priors generated from the active specifications; positive wide-scale parameters use labelled natural-log plotting coordinates, including the correct prior Jacobian. Broad priors are not renormalized to the displayed posterior window; full support is in the prior table. In each two-parameter marginal, KDE contours enclose approximately 68\% and 95\% probability. Interfacial energy is fixed at 0.64 mJ per squared metre and is absent from these inferred coordinates. Orange markers denote the transformed inference-coordinate MAP. Logarithmic plotting coordinates use numerical values of AJ in particles/(L min), AG in micrometres/min, and D0 in squared micrometres/min; no reference-rate parameters are introduced. Sampling convergence does not establish numerical convergence, identifiability or physical model adequacy.}
\newcommand{\ProteinCNTFitCaption}{CNT nucleation and size-independent mean growth. \ProteinFitDetails MUSCL growth and native SSPRK33 use 400 bins on 0--50 $\mu$m with combined CFL 0.35 and maximum step 2 min; 200 retained parameter draws supply the bands. The posterior retains 4 chains of 750 draws each.}
\newcommand{\ProteinCNTCornerCaption}{CNT nucleation and size-independent mean growth: NUTS posterior. \ProteinCornerDetails}
\newcommand{\ProteinGRDFitCaption}{CNT nucleation and mean growth with constant random-fluctuation diffusivity $D_0$. \ProteinFitDetails MUSCL growth and native SSPRK33 use 200 bins on 0--50 $\mu$m with combined CFL 0.35 and maximum step 2 min; 200 retained parameter draws supply the bands. The posterior retains 4 chains of 750 draws each.}
\newcommand{\ProteinGRDCornerCaption}{CNT nucleation and mean growth with constant random-fluctuation diffusivity $D_0$: NUTS posterior. \ProteinCornerDetails}
\newcommand{\ProteinPriorTableCaption}{Implemented priors, physical units and applicable models. Inference transformations are separate from prior-density coordinates; the dispersion Exponential prior has mean/scale 0.01 $\mu{\rm m}^2/{\rm min}$.}
\newcommand{\ProteinParameterTableCaption}{Inference-coordinate MAP estimates transformed to physical units, posterior medians and equal-tailed 95\% credible intervals. Model 1 fixes dispersion to zero. Intervals condition on the model, fixed observation errors and physical constants.}
